% !PRINTED 639
\weberpage{639}
%<<CONT
1. \gesperrt{Die einzigen in $\Omega_m$ enthaltenen Einheitswurzeln
sind die mit positivem und negativem Zeichen
genommenen Potenzen von $r$.}

Hieran schliesst sich ein ganz allgemeiner, d.\,h. in jedem
algebraischen Zahlkörper gültiger Satz über die Einheitswurzeln,
den wir jetzt beweisen wollen.

Der Satz lautet:

2. \gesperrt{Ist $\alpha$ eine ganze Zahl eines algebraischen Zahlkörpers
 $n^{\text{ten}}$ Grades, und haben alle mit $\alpha$ conjugirten
 Zahlen $\alpha_1, \alpha_2, \dots, \alpha_n$ den absoluten
Werth 1, so ist $\alpha$ eine Einheitswurzel}\footnote{\gesperrt{Kronecker}, \glqq Zwei Sätze über Gleichungen mit ganzzahligen
Co\"efficienten\grqq. Crelle's Journal, Bd. 53 (1857). \gesperrt{Minkowski}, \glqq Geometrie
der Zahlen\grqq, Art. 43.}.

Zum Beweis ist zunächst zu bemerken, dass der absolute
Werth der Norm einer ganzen Zahl $\omega$ dem Product der absoluten
Werthe der conjugirten Zahlen $\omega_1, \omega_2, \dots, \omega_n$ gleich, und als
ganze rationale Zahl jedenfalls grösser oder gleich 1 ist. Es ist
daher nicht möglich, dass alle diese absoluten Werthe kleiner
als 1 sind. Sind sie aber alle gleich 1, wie bei der in unserem
Satze angenommenen Zahl $\alpha$, so muss die Norm $= \pm 1$ sein,
und $\alpha$ ist eine Einheit. Ist $\beta$ eine zweite Einheit, die mit ihren
conjugirten zugleich den absoluten Werth 1 hat, so hat der
Quotient $\alpha : \beta$ dieselbe Eigenschaft. Wenn daher unter den zu
diesem Quotienten conjugirten Zahlen auch nur eine reell ist,
so muss
\[ \frac{\alpha}{\beta} = \pm 1 \]
oder $\alpha = \pm \beta$ sein, und dies gilt dann auch noch, wenn $\alpha, \beta$
durch die entsprechenden Zahlen eines conjugirten Körpers
ersetzt werden.

Wenn also $\alpha$ und $\beta$ weder gleich noch entgegengesetzt sind,
so ist ihr Verhältniss nicht reell, und es besteht daher zwischen
den absoluten Werthen die Ungleichung:
\[ |\alpha \pm \beta| < |\alpha| + |\beta|. \]
(Vergl. die Einleitung zum ersten Bande, S. 19.)

Es ist also
\[ |\alpha \pm \beta| < 2, \]
%CONT>>
